\begin{table}[H]
\caption{Instruction Decoder Service Requests}
\label{tab:asDecServiceReq}
\centering
\begin{tabularx}{\textwidth}{|l |X|}
  \hline
  Service Request & Explanation \\
  \hline
  \hline
  \texttt{trap\_illegal\_instr\_o} &
    Synchronous, combinational request: any \texttt{instr\_opcode\_i}/
    \texttt{instr\_func3\_i} combination that does not decode to a valid
    RV64I instruction. This is the \emph{sole} source of the
    illegal-instruction exception (\texttt{mcause}\,=\,2) on both chip
    variants -- the decoder itself does not distinguish \texttt{ecall},
    \texttt{ebreak}, or \texttt{mret} (all legal opcode=115 encodings from
    its point of view); those are recognized directly on the instruction
    word by the CPU top level. See Section~\ref{subsec:irqconcept}
    (Interrupt Concept) for the full trap-priority chain and CSR update
    sequence this signal feeds into. \\
  \hline
\end{tabularx}
\end{table}

The Instruction Decoder has no other service request: it is a
combinational block with no clock, no interrupt/DMA request of its own
kind, and no error condition beyond the single illegal-instruction case
above.
